% ===========================================================================
%  15 — Engenharia de entrega
% ===========================================================================
\chapter{Engenharia de entrega}
\label{cap:entrega}

\section{O \emph{monorepo} e o \emph{build} incremental}

\begin{decisao}
\textbf{\emph{Monorepo} único, com pnpm \emph{workspaces} + Turborepo}
(\adrr{ADR-02}). A SDK é consumida por 48 aplicações: uma mudança nela precisa
testar tudo no \emph{mesmo} \emph{pull request}, ou a incompatibilidade aparece
num \emph{bundle} publicado semanas depois. Repositórios separados só se houver
estúdio externo produzindo experiência.
\end{decisao}

\begin{lstlisting}[language=jsonunoo, caption={\arq{turbo.json} --- o \emph{pipeline} inteiro}, label={lst:turbo}]
{
  "tasks": {
    "build":     { "dependsOn": ["^build"],
                   "outputs": ["dist/**", ".next/**", "!.next/cache/**"] },
    "dev":       { "cache": false, "persistent": true },
    "lint":      { },
    "typecheck": { "dependsOn": ["^build"] },
    "test":      { "dependsOn": ["^build"] }
  }
}
\end{lstlisting}

O \cd{dependsOn: ["\^{}build"]} é o que faz o grafo funcionar: \cd{typecheck} de
uma experiência exige o \cd{build} da SDK, então o \emph{filtro} por mudança
arrasta as dependências corretas.

\begin{tabularx}{\linewidth}{@{}L{50mm}X@{}}
\toprule
\headrow \thd{Comando} & \thd{Uso}\\
\midrule
\cd{pnpm turbo build --filter=@edu/exp-01} & Builda uma experiência mais as
  dependências dela.\\
\zebra \cd{pnpm turbo build --filter='...[HEAD\^{}]'} & Só o que mudou. É o que o
  CI usa.\\
\cd{pnpm turbo typecheck lint test} & O portão local da \emph{Definition of
  Done}.\\
\bottomrule
\end{tabularx}

\section{Integração contínua}

\begin{lstlisting}[language=yamlunoo, caption={\arq{.github/workflows/ci.yml} (abreviado)}, label={lst:ci}]
name: ci
on:
  # Todo PR: de branch de trabalho para dev-local, dai para dev, para hml
  # e, quando producao existir, para main — o mesmo portao nas quatro pontes.
  pull_request:
  # Chamado pelo deploy.yml no push para dev/hml. NAO tem `push` proprio:
  # rodaria duas vezes no mesmo commit, e o deploy precisa ESPERAR por ele —
  # dois workflows em paralelo deixariam um teste vermelho subir para escola.
  workflow_call:
    inputs:
      completo:   # `forcar` no deploy builda tudo; o portao cobre o mesmo
        type: boolean

jobs:
  verify:
    runs-on: ubuntu-latest
    services:
      postgres:   # 9 das 23 suites da api sao de INTEGRACAO (RLS, locks)
        image: postgres:16
    steps:
      - uses: actions/checkout@v4
        with: { fetch-depth: 2 }
      - uses: pnpm/action-setup@v4
      - uses: actions/setup-node@v4
        with: { node-version: 24, cache: pnpm }
      - run: pnpm install --frozen-lockfile
      - run: pnpm --filter @edu/api migrate:deploy
      - run: |
          if [ "$COMPLETO" = "true" ]; then
            pnpm turbo typecheck lint test build
          else
            pnpm turbo typecheck lint test build --filter='...[HEAD^]'
          fi
      # COMENTADO no arquivo real desde 07d9ec0 -- ver a advertencia abaixo
      - name: Snyk (trava em high)   # high/critical em dependencia = vermelho
        env: { SNYK_TOKEN: "${{ secrets.SNYK_TOKEN }}" }
        run: |   # o --exclude e obrigatorio: o passo roda DEPOIS do build
          npm install -g snyk
          snyk test --all-projects --dev --severity-threshold=high \
            --exclude=.next,.next-prod,dist,build,storage,storage-privado,jogo
\end{lstlisting}

\begin{decisao}
O CI é chamado pelo \emph{deploy} por \cd{workflow\_call}, e \textbf{não} tem
gatilho próprio em \cd{push}. Com gatilho próprio, ele rodaria duas vezes
no mesmo \emph{commit} --- e, pior, em \emph{paralelo} com o \emph{deploy}, o que
deixaria um teste vermelho subir para uma escola. Assim o \emph{deploy}
\textbf{espera} pelo portão. O gatilho de \cd{pull\_request} não filtra a
\emph{branch} de destino de propósito: o portão é o mesmo nas três pontes do
fluxo (\emph{branch} de trabalho \(\to\) \cd{dev-local}, \cd{dev-local}
\(\to\) \cd{dev}, \cd{dev} \(\to\) \cd{hml}, \cd{hml} \(\to\) \cd{main}).
\end{decisao}

\subsection{A trava de vulnerabilidade (Snyk)}

\begin{atencao}
\textbf{O passo está COMENTADO no \arq{ci.yml} desde o \emph{commit}
\cd{07d9ec0}} (\enquote{\emph{chore: disable Snyk vulnerability scan in CI
workflow}}), situação de 2026-09-05. O parágrafo abaixo descreve o portão como
ele foi desenhado e como volta a funcionar quando o passo for reativado ---
\textbf{hoje ele não trava nada}. O bloco de YAML acima mostra o passo na forma
ativa, que é a forma comentada no arquivo. Ver a advertência da \adrr{ADR-39},
inclusive para o efeito colateral: o prazo do \arq{apps/api/.snyk} era cobrado
por este passo e agora vence em silêncio.
\end{atencao}

O portão termina com
\cd{snyk test --all-projects --dev --severity-threshold=high}: vulnerabilidade
\cd{high} ou \cd{critical} em qualquer dependência do \emph{workspace} deixa o
CI vermelho --- e, como o \emph{deploy} espera pelo portão, nada sobe
(\adrr{ADR-39}). O passo roda sempre no \emph{workspace} inteiro, sem o filtro
por mudança do resto do portão: o \emph{lockfile} é um só, e a vulnerabilidade
nova chega pelo banco de dados do Snyk, não por um \emph{commit} ---
``dependência que não mudou'' não existe como categoria.

O \cd{--dev} não é detalhe: sem ele o Snyk pula as \cd{devDependencies}, e foi
por aí que um \emph{critical} no \cd{vitest} passou verde na primeira semana da
trava (achado pelo \cd{pnpm audit}, que olha tudo). Ferramenta de
desenvolvimento não embarca no \emph{bundle}, mas roda na máquina de quem
desenvolve e no \emph{runner} do CI --- e é lá que um \emph{dev server} furado
vaza código.

\begin{nota}
\textbf{O que exige \emph{major} de \emph{framework} não bloqueia --- expira.}
O \arq{apps/api/.snyk} ignora, com \textbf{prazo e motivo}, as vulnerabilidades
cujo \emph{fix} é o Nest~11 (\cd{@nestjs/core}, e o \cd{file-type} que é ESM-only
e quebraria o \cd{require} do Nest~10) ou um \emph{bump} que só o Prisma pode
fazer (\cd{deepmerge-ts}). Quando o prazo vence sem o \emph{upgrade}, o próprio
passo volta a falhar --- é o prazo \textbf{com} a rotina que o cumpre, a
resposta à armadilha do \cd{expires\_at}. O resto foi consertado na hora: o hub
subiu para o \textbf{Next~16} (as \emph{criticals} do 15 não tinham \emph{patch}
na linha 15.x), e \cd{pnpm.overrides} no \arq{package.json} da raiz cobre o
transitivo (\cd{next} 16.3.3 por baixo do \emph{design system}, \cd{multer}
2.2.0, \cd{js-yaml} 4.3.1, \cd{lodash} 4.18.1).
\end{nota}

O segredo \cd{SNYK\_TOKEN} vem dos \emph{secrets} do repositório
(\cd{snyk config get api | gh secret set SNYK\_TOKEN}). Ausente, o passo
\textbf{falha alto} --- nunca pula, pela mesma regra do \emph{guard} do
Pub/Sub: env ausente é porta fechada, não porta aberta.

\section{Entrega contínua}

\begin{figure}[H]
\centering
\fitpage{%
\begin{tikzpicture}[x=1cm, y=1cm]

  \node[etapaio, text width=26mm] (push) at (1.4,0)
    {\emph{push} para \texttt{dev} ou \texttt{hml}\desc{merge de um PR}};
  \node[etapa, text width=26mm, fill=unooGreenWash, draw=unooGreen] (ver) at (5.0,0)
    {\textbf{portão}: \texttt{ci.yml}\desc{typecheck, lint, test, build}};
  \node[etapa, text width=30mm] (diff) at (9.0,0)
    {\emph{job} \texttt{mudancas}\desc{\texttt{git diff HEAD\^{} HEAD}\\api $|$ hub $|$ exps}};

  % --- trilha da api (job proprio)
  \node[camadalbl, anchor=west] at (0.4,-3.6)
    {\emph{job} \texttt{deploy-api} --- a ORDEM importa};
  \node[etapa, text width=28mm] (img) at (2.4,-4.8)
    {imagem da api\desc{camada de deps do \texttt{pnpm fetch}, \emph{cache} GHA}};
  \node[etapa, text width=30mm, fill=unooYellowWash, draw=unooYellow] (mig) at (6.4,-4.8)
    {\textbf{Cloud Run Job} \texttt{-migrate}\desc{na VPC, \texttt{--wait}}};
  \node[etapaout, text width=28mm] (dep) at (10.4,-4.8)
    {\texttt{gcloud run deploy}\desc{revisão nova}};

  % --- trilha do hub (job proprio, em paralelo com a api)
  \node[camadalbl, anchor=west] at (0.4,-6.4)
    {\emph{job} \texttt{deploy-hub} --- em PARALELO com a api};
  \node[etapa, text width=28mm] (imgh) at (2.4,-7.4)
    {imagem do hub\desc{camada de deps do \texttt{pnpm fetch}, \emph{cache} GHA}};
  \node[etapaout, text width=28mm] (deph) at (6.4,-7.4)
    {\texttt{gcloud run deploy}};

  % --- trilha das experiencias (dois jobs: subir e ativar)
  \node[camadalbl, anchor=west] at (0.4,-10.6)
    {\emph{jobs} \texttt{publicar-exps} (em paralelo) \seta{} \texttt{ativar-exps} (depois dos dois)};
  \node[etapa, text width=24mm] (build) at (1.8,-11.8)
    {\texttt{turbo build}\desc{filtrado, ou tudo se \texttt{forcar}}};
  \node[etapa, text width=30mm] (rsync) at (5.4,-11.8)
    {\texttt{gcloud storage rsync}\desc{\texttt{gs://\{bucket\}/}\\\texttt{\{exp\}/\{ver\}/}}};
  \node[etapa, text width=28mm] (cc) at (9.2,-11.8)
    {\texttt{Cache-Control}\desc{\texttt{max-age=60} só no \texttt{index.html}}};
  \node[etapaout, text width=30mm] (man) at (13.0,-11.8)
    {\texttt{publicar-manifesto.mjs}\desc{\emph{job} \texttt{ativar-exps}:
     \texttt{PUT /experiencias/:expId}, token de 5 min}};

  \draw[pipe] (push) -- (ver);
  \draw[pipe] (ver) -- (diff);
  \draw[pipe] (diff.south) -- (9.0,-1.6) -- (-0.4,-1.6) -- (-0.4,-4.8) -- (img.west);
  \draw[pipe] (-0.4,-4.8) -- (-0.4,-7.4) -- (imgh.west);
  \draw[pipe] (img) -- (mig);
  \draw[pipe] (mig) -- (dep);
  \draw[pipe] (imgh) -- (deph);
  \draw[pipe] (dep.south) -- (10.4,-9.4) -- (6.4,-9.4);
  \draw[pipe] (deph.south) -- (6.4,-9.4) -- (-0.4,-9.4) -- (-0.4,-11.8) -- (build.west);
  \draw[pipe] (build) -- (rsync);
  \draw[pipe] (rsync) -- (cc);
  \draw[pipe] (cc) -- (man);

  \node[actnote, text width=44mm, anchor=north west] at (12.2,-3.4)
    {\textbf{Por que a migração antes do \emph{deploy}:} o \emph{schema} novo tem
     de estar no banco \emph{antes} de a revisão nova receber tráfego, senão a
     API roda contra um banco velho. O \texttt{--wait} faz o passo falhar se a
     migração falhar --- e o \emph{deploy} não acontece.};

  \node[actnote, text width=44mm, anchor=north west] at (12.2,-6.9)
    {\textbf{Por que as experiências depois do hub:} o \texttt{PUT} do manifesto é
     validado pela API contra o JWKS que o \emph{hub} serve. Hub fora do ar =
     manifesto rejeitado --- e um manifesto rejeitado é melhor que um manifesto
     aceito apontando para um hub que não subiu.};

\end{tikzpicture}}
\caption[Pipeline de entrega]{O \emph{pipeline} de entrega. O portão e o
  \emph{job} \cd{mudancas} decidem o que roda; api e hub sobem em \emph{jobs}
  paralelos (cada um autentica por WIF, sem chave JSON); as experiências
  publicam por último, depois dos dois.}
\label{fig:pipeline-deploy}
\end{figure}

\subsection{Sonda, \emph{smoke} e aviso (2026-09-02)}

Três passos que o \emph{pipeline} não tinha, todos da fase 0 do
Capítulo~\ref{cap:observabilidade} --- a parte dela que é deste repositório:

\begin{itemize}
  \item \textbf{Sonda pós-\emph{deploy}} nos \emph{jobs} da api e do hub. O
    \cd{gcloud run deploy} volta quando a revisão \emph{recebe tráfego}, não
    quando ela responde; o passo pergunta, pelo domínio público, até 12 vezes
    em 2 minutos: \rota{/readyz} \(=\) \cd{\{"ok":true\}} na api (processo +
    banco), e \rota{/api/healthz} \textbf{e} \rota{/.well-known/jwks.json}
    com \cd{"keys"} no hub --- sem a chave carregada nenhum \emph{token} de
    experiência é validável, e o \cd{ativar-exps} falharia com um 401 que não
    aponta para a causa. Vermelho = o \cd{ativar-exps} não ativa nada.
  \item \textbf{\emph{Smoke}}: o \emph{job} \cd{smoke} roda a suíte da
    \emph{landing} (\cd{publico/01}) de \arq{scripts/e2e} contra o ambiente
    recém-publicado, depois de tudo --- só leitura e sem \emph{login}, então sem
    senha de demonstração no CI; o relatório sobe como artefato em falha.
    Falha de \emph{workflow} avisa por e-mail quem fez o \emph{push}
    (comportamento padrão do GitHub); o alerta de disponibilidade em regime é
    dos \emph{uptime checks} do Cloud Monitoring (Terraform), por e-mail. A
    suíte agendada como monitor sintético (a \enquote{sonda} da
    Seção~\ref{sec:ferramentas-obs}) foi deixada para uma onda posterior, por
    decisão de 2026-09-02.
\end{itemize}

E um quarto passo, dos \emph{sourcemaps}: o \cd{rsync} das experiências ganhou
\cd{-{}-exclude='.*\textbackslash.map\$'} (vale para origem e destino), e os
\cd{.map} sobem para o \emph{bucket} \textbf{privado} \env{EXP\_MAPAS\_BUCKET},
no mesmo caminho versionado. \arq{scripts/mapear-stack.mjs} traduz a
\emph{stack} de um erro de cliente com o mapa da versão exata.

\subsection{O que faz o \emph{deploy} caber no café}

\begin{decisao}
\textbf{As imagens de api e hub sobem em \emph{jobs} paralelos, e a camada de
dependências dos \emph{Dockerfiles} é cacheável por si} (\adrr{ADR-38}). Medido
nos \emph{runs} de 2026-08-28: as duas imagens eram 80\% do tempo de
\emph{deploy} --- 4m53s a api num \emph{push} comum, mais 6m30s o hub em série
num \cd{forcar} --- porque \cd{COPY . .} vinha \emph{antes} do
\cd{pnpm install}, e qualquer mudança de código invalidava a camada inteira. O
\emph{cache} GHA que o \emph{workflow} sempre teve só acertava apt e corepack.
\end{decisao}

Os dois \emph{Dockerfiles} passaram a abrir com uma camada que depende
\textbf{só do \emph{lockfile}}:

\begin{lstlisting}[language=bashunoo, caption={A camada de dependências (api e hub)}]
# Camada de dependencias: so invalida quando o lockfile muda.
COPY pnpm-lock.yaml pnpm-workspace.yaml ./
RUN pnpm fetch
# Codigo por cima; o install resolve do store local, sem rede.
COPY . .
RUN pnpm install --frozen-lockfile --offline
\end{lstlisting}

\cd{pnpm fetch} existe para isto: baixa o \emph{store} inteiro lendo apenas o
\emph{lockfile}, sem precisar dos \arq{package.json}. Num \emph{commit} que só
toca código, a camada cara vem do \emph{cache} e sobra o que é daquele
\emph{app}: \cd{nest build} (ou \cd{next build}), \cd{pnpm deploy --prod} e o
\cd{prisma generate}.

E como as duas imagens não dependem uma da outra, cada uma virou um \emph{job}
--- \cd{deploy-api} e \cd{deploy-hub} --- atrás do mesmo portão, com o
\emph{diff} calculado uma única vez num \emph{job} \cd{mudancas}. A ordem que
\textbf{importa} continua explícita no grafo, agora em duas partes: SUBIR o
\emph{bundle} não depende de nada além do portão, porque o caminho carrega a
versão (\cd{\{expId\}/\{versao\}/}) e ninguém pede por ele até o manifesto
mudar --- o \emph{upload} é inerte. ATIVAR depende dos dois \emph{deploys}, e
isso não é preferência: a api valida o \emph{PUT} contra o JWKS que o hub serve.
Separar os dois encurta a janela em que o hub novo conviveria com o
\emph{bundle} antigo: de um \cd{turbo build} inteiro para a duração de um
\emph{PUT}. E dentro de \cd{deploy-api} a migração continua antes da revisão
nova. A única ordem que se
perdeu foi acidental: o hub não espera mais a api --- os dois conversam por
HTTP e o \emph{rollout} do Cloud Run nunca garantiu revisão casada mesmo.

\subsection{Nada de nome de recurso no \emph{workflow}}

\begin{decisao}
Tudo que muda de um ambiente para outro entra por \cd{vars} de
\textbf{\emph{environment}} do GitHub (\cd{dev}, \cd{hml} e, no futuro,
\cd{prd}) --- o \emph{workflow} \textbf{não sabe} o nome de nenhum recurso. Cada
\emph{job} de \emph{deploy} declara \cd{environment:} com o nome da
\emph{branch} que disparou o \emph{push} (\cd{main} mapeia para \cd{prd}), e é
esse vínculo que resolve as \cd{vars} da \emph{stack} certa. Os valores saem de
\cd{terraform output} do \emph{workspace} correspondente, e há um \emph{script}
que os grava no \emph{environment} certo
(\arq{scripts/configurar-github-ci.sh}) --- ver a \adrr{ADR-53}.

É isso que faz o mesmo \emph{workflow} servir a \cd{dev}, a \cd{hml} e à
produção, e ao próximo contrato, sem \emph{fork}.
\end{decisao}

\begin{atencao}
As \cd{vars} de \emph{environment} \textbf{só resolvem dentro de um \emph{job}
que declara \cd{environment:}} --- um bloco \cd{env:} no nível do
\emph{workflow} as receberia vazias, sem erro. Por isso o \arq{deploy.yml} não
tem \cd{env:} global: cada \emph{job} de \emph{deploy} declara o
\emph{environment} e as próprias variáveis. E referenciar um \emph{environment}
que não existe o \textbf{cria em silêncio}, com todas as \cd{vars} vazias ---
o \emph{job} \cd{mudancas} tem um passo de guarda que derruba um
\cd{workflow\_dispatch} disparado de \emph{branch} que não é de ambiente.
\end{atencao}

\begin{tabularx}{\linewidth}{@{}L{34mm}X@{}}
\toprule
\headrow \thd{\cd{vars}} & \thd{Conteúdo}\\
\midrule
\cd{GCP\_PROJECT}, \cd{REGION} & Projeto e região.\\
\zebra \cd{AR\_REPO} & Artifact Registry.\\
\cd{API\_SERVICE}, \cd{HUB\_SERVICE}, \cd{MIGRATE\_JOB} & Nomes dos serviços e do
  \emph{job}.\\
\zebra \cd{EXP\_BUCKET}, \cd{ASSETS\_BASE\_URL} & \emph{Buckets} e base pública de
  \emph{assets}.\\
\cd{API\_URL}, \cd{HUB\_ORIGIN} & Endereços públicos.\\
\zebra \cd{WIF\_PROVIDER}, \cd{DEPLOY\_SA} & Identidade federada do CI.\\
\cd{CLARITY\_PROJECT\_ID} & Id do projeto Microsoft Clarity
  (\adrr{ADR-60}). Única \cd{vars} que \textbf{não} sai de
  \cd{terraform output}: é criada à mão em cada \emph{environment}.\\
\bottomrule
\end{tabularx}

\begin{atencao}
\textbf{\cd{CLARITY\_PROJECT\_ID} entra na \emph{imagem}, não no serviço.} É a
única variável de todo o \emph{deploy} que chega ao \emph{bundle} do
\textbf{cliente}, e \cd{NEXT\_PUBLIC\_*} é substituída pelo literal durante o
\cd{next build} --- defini-la no Cloud Run depois não altera um \emph{bundle}
já compilado. Por isso o passo \emph{Imagem do hub} a passa como
\cd{build-arg}, e o \arq{apps/hub/Dockerfile} a recebe num \cd{ARG} antes do
\cd{turbo build}. As duas consequências: trocar o id exige \emph{rebuild} (não
basta reiniciar o serviço), e uma imagem passa a ser específica do ambiente que
a construiu --- promover \cd{hml} para produção reconstrói, não repromove a
mesma \emph{tag}. Ausente, o \cd{ARG} cai no \emph{default} vazio e o hub sobe
com o Clarity desligado, sem erro.
\end{atencao}

\subsection{Concorrência e reexecução}

\begin{lstlisting}[language=yamlunoo, caption={A política de concorrência do \emph{deploy}}]
# Dois deploys ao mesmo tempo brigam pela mesma revisao do Cloud Run e pelo
# mesmo banco. `cancel-in-progress: false` de proposito: abortar no meio de
# uma migration deixa schema e imagem fora de sincronia — melhor enfileirar.
concurrency:
  group: deploy-${{ github.ref }}
  cancel-in-progress: false
\end{lstlisting}

\begin{nota}
Há um \cd{workflow\_dispatch} com \emph{input} \cd{forcar}, para publicar tudo
ignorando o \emph{diff} --- primeiro \emph{deploy} de um contrato, ou reexecução
depois de falha. Com \cd{forcar}, o \emph{build} também tem de ser de
\textbf{todas} as experiências, senão o laço de publicação encontra pasta sem
\arq{dist/}.
\end{nota}

\begin{nota}
O laço de publicação é tolerante por desenho: pasta de experiência sem
\arq{package.json}, sem \arq{public/experiencia.json} ou sem \arq{dist/} é
\textbf{pulada com aviso}, não é erro. Experiência incompleta é trabalho em
andamento, e trabalho em andamento não deve quebrar o \emph{deploy} do resto.
\end{nota}

\section{Ambientes locais}

\begin{tabularx}{\linewidth}{@{}L{30mm}L{34mm}X@{}}
\toprule
\headrow \thd{Modo} & \thd{Comando} & \thd{Quando usar}\\
\midrule
Banco em contêiner, resto no \emph{host} &
  \cd{docker compose up -d pg} + \cd{pnpm dev} &
  Iteração rápida com \emph{debugger} do editor. O caminho mais direto.\\
\zebra \textbf{Tudo em contêiner, sincronizado} &
  \cd{pnpm docker:dev} &
  \textbf{Recomendado.} Banco, API, hub e as seis experiências, cada um com o
  próprio \emph{watcher}, com o código do \emph{host} por \emph{bind mount}:
  salvar um arquivo recompila lá dentro (API $\sim$1\,s, hub e Vite na hora), sem
  \cd{build} e sem \cd{up} de novo.\\
Imagens congeladas &
  \cd{docker compose up -d --build} &
  Aproximar-se da produção. Migrações e \emph{seed} rodam sozinhos no primeiro
  \cd{up}.\\
\bottomrule
\end{tabularx}

\begin{atencao}
As portas são as mesmas nos três modos (hub \cd{:3000}, API \cd{:3001},
experiências \cd{:5173}--\cd{:5177}), então \textbf{não dá para rodar
\cd{pnpm docker:dev} e \cd{pnpm dev} ao mesmo tempo} --- é um \emph{ou} o outro.
\end{atencao}

\begin{tabularx}{\linewidth}{@{}L{40mm}X@{}}
\toprule
\headrow \thd{Situação} & \thd{O que fazer}\\
\midrule
Mexeu em dependência (\arq{package.json} ou \emph{lockfile}) &
  \cd{pnpm docker:dev:deps} e reiniciar. O \arq{node\_modules} de cada contêiner
  é \textbf{volume próprio} --- binário de macOS não roda no Linux da imagem ---
  então instalar no \emph{host} não basta.\\
\zebra Salvou um arquivo e nada recompilou &
  \cd{DEV\_POLLING=1 pnpm docker:dev}. O \emph{inotify} pode não atravessar o
  \emph{bind mount} em alguns \emph{setups}.\\
\bottomrule
\end{tabularx}

\section{O provedor de identidade sobe à parte}

\begin{lstlisting}[language=bashunoo, caption={O WSO2 tem ciclo de vida próprio}]
docker compose -f wso2/docker-compose.yml up -d     # IS :9443, banco dele :5433
docker compose -f wso2/docker-compose.yml watch     # olha conf/ e bootstrap/
\end{lstlisting}

\begin{nota}
O \cd{watch} é o equivalente do \emph{hot reload} para um servidor Java:
\arq{conf/} já entra por volume (nada de \emph{rebuild}), e mudar o
\arq{deployment.toml} \textbf{reinicia} o servidor --- que é o ciclo mais curto
possível ali, porque o \emph{entrypoint} só copia o volume de configuração para o
\emph{server home} no \emph{boot} ($\sim$2\,min). Mudar \arq{bootstrap/}
reexecuta o \emph{bootstrap} (idempotente); só mudar o \arq{Dockerfile} (o
\emph{driver} JDBC) refaz imagem.
\end{nota}

\section{Convenções de \emph{git}}

\begin{tabularx}{\linewidth}{@{}L{34mm}X@{}}
\toprule
\headrow \thd{Convenção} & \thd{Regra}\\
\midrule
Fluxo & \textbf{Quatro ambientes} (\adrr{ADR-101}): \emph{branch} de trabalho
  \(\to\) PR para \arq{dev-local} \(\to\) PR para \arq{dev} \(\to\) PR para
  \arq{hml} \(\to\) (futuro) PR para \arq{main}. Quem testa muda a cada degrau:
  \arq{dev-local} é o Docker da máquina e quem olha é quem desenvolve;
  \arq{dev} é do time; \arq{hml}, do cliente. \textbf{Ninguém commita direto em
  \emph{branch} de ambiente} --- em \arq{dev}, \arq{hml} e \arq{main} o
  \emph{merge} do PR dispara o \emph{deploy} da \emph{stack}; em
  \arq{dev-local} ele integra a frente ao ambiente local, que não tem
  \emph{deploy} e por isso não está no \arq{deploy.yml}.\\
\zebra \emph{Worktree} & \textbf{Toda tarefa começa num}, paralela ou não
  (\adrr{ADR-61}, \adrr{ADR-101}): \cd{pnpm worktree criar <branch> --fatia
  <caminhos>}, sobre \cd{origin/dev-local} ESCRITO. A cópia principal é onde a
  \emph{stack} roda, e editar nela mistura o que subiu com o que está sendo
  escrito. \cd{--fatia} é obrigatória (\adrr{ADR-71}).\\
\emph{Merge} & \emph{Merge commit} ou \emph{squash}; \textbf{rebase fica
  desabilitado} nas opções do repositório. O \emph{job} \cd{mudancas} compara
  \cd{HEAD\^{}} (o primeiro pai), que num \emph{merge commit} é a ponta antiga
  da \emph{branch} de ambiente --- um \emph{rebase-merge} faria o \emph{diff}
  enxergar só o último \emph{commit} do PR.\\
\zebra \emph{Branches} & \textbf{Uma por experiência}, com o nome do slug
  (\arq{exp-12-fracoes}). Trabalho que não é experiência usa o prefixo do tipo do
  \emph{commit}: \arq{fix/quota-por-perfil}, \arq{chore/deps},
  \arq{ci/deploy}.\\
\emph{Commits} & \cd{feat(exp-12): ...}, \cd{fix(api): ...},
  \cd{chore(sdk): ...} --- o escopo é o slug da experiência ou o aplicativo
  tocado. Validado pelo \emph{hook} \arq{.githooks/commit-msg}.\\
\zebra Promoção & O PR de \cd{hml} para \cd{main} leva \textbf{tag semver}
  (\cd{v1.4.0}) e \emph{release}, criadas depois do \emph{merge}. MAJOR só em
  quebra de contrato externo (\cd{PROTOCOL\_VERSION} da ponte, manifesto,
  telemetria, contrato do Terraform); MINOR é o caso comum. A \emph{tag} é o
  nome do que está em produção --- sem ela, \enquote{voltar para a anterior}
  vira caça a \emph{SHA} no meio de um incidente.\\
\emph{Hotfix} & Sai de \cd{main}, entra por PR em \cd{main}, leva
  \emph{tag} de PATCH --- e desce a escada INTEIRA por \textbf{retro-merge} no
  mesmo dia: \cd{hml}, \cd{dev} e \cd{dev-local}. Parar em \cd{dev} deixa a
  base local sem a correção, e a frente seguinte nasce sem ela.\\
\zebra Imposição & \cd{branch protection} e \cd{rulesets} respondem \textbf{403}
  (repositório privado, plano \emph{Free}). As regras vivem em \arq{.githooks/}
  --- \arq{pre-push} recusa \emph{push} direto em \emph{branch} de ambiente e
  nome fora do padrão --- instaladas pelo \cd{prepare} do \cd{pnpm install}
  (\cd{core.hooksPath}). Ver \adrr{ADR-56}.\\
\bottomrule
\end{tabularx}

\begin{nota}
O fluxo completo, com os comandos de cada ponte (promoção, \emph{hotfix},
\emph{retro-merge}, \emph{tag}, \emph{rollback} por camada), está em
\arq{CONTRIBUTING.md}, na raiz do repositório.
\end{nota}

\begin{atencao}
\textbf{Nada de \emph{build} no \emph{git}.} O \arq{.gitignore} cobre
\arq{dist/}, \arq{.next/}, \arq{.next-prod/}, \arq{.turbo/},
\arq{*.tsbuildinfo}, \arq{storage/}, \arq{storage-privado/} e os \emph{prints}
das suítes de ponta a ponta.

\arq{apps/hub/.next-prod/} já entrou uma vez por um \cd{git add -A} distraído:
358 arquivos e \SI{250}{\mega\byte} de \emph{cache} de \emph{webpack} no
histórico, que só saíram reconstruindo a \emph{branch} inteira.
\textbf{Stage por caminho, nunca \cd{-A} sem olhar o \cd{git status}.}
\end{atencao}

\begin{nota}
\arq{apps/hub/next-env.d.ts} \textbf{é} versionado (o Next espera isso), mas
alterna entre \arq{./.next} e \arq{./.next-prod} conforme o último \emph{build}.
Modificado depois de um \cd{build:isolado} = descartar, não \emph{commitar}.
\end{nota}
